\section{Statistical Models}

This section presents the seven statistical models developed to understand African Swine Fever (ASF) transmission risk in Romania. We employed three distinct variable selection philosophies---theory-driven, information-theoretic, and regularization-based---to triangulate robust predictors of ASF occurrence. Additionally, we developed domain-specific models to assess the relative importance of different ecological factors.

\subsection{Overview of Modeling Strategy}

Our modeling framework addresses a fundamental question in disease ecology: \textit{Which environmental and epidemiological factors best predict ASF transmission to domestic pigs?} We approached this question from multiple angles:

\begin{enumerate}
    \item \textbf{Theory-Driven Model}: Based on established ecological hypotheses about ASF transmission
    \item \textbf{Stepwise AIC Selection}: Data-driven selection balancing fit and parsimony
    \item \textbf{Elastic Net Regularization}: Machine learning approach optimizing predictive performance
    \item \textbf{Domain-Specific Models}: Four separate models isolating effects of Livestock, Landscape, Outbreak History, and Wildlife factors
\end{enumerate}

This triangulation strategy allows us to identify robust predictors that emerge across different selection methods, while also understanding which variables are selected by specific approaches and why.

%======================================================================
\subsection{Model 1: Theory-Driven Model (11 Variables)}
%======================================================================

\subsubsection{Rationale and Variable Selection}

The theory-driven model was constructed \textit{a priori} based on established ecological principles of ASF transmission. Rather than letting data guide variable selection, we specified 11 variables that represent key biological mechanisms through which ASF spreads among domestic pigs.

\begin{table}[H]
\centering
\caption{Theory-Driven Model Variables: Ecological Hypotheses and Biological Rationale}
\label{tab:theory_driven_variables}
\small
\begin{tabular}{p{3.5cm}p{2cm}p{2.5cm}p{6cm}}
\toprule
\textbf{Variable} & \textbf{Domain} & \textbf{Hypothesis} & \textbf{Biological Rationale} \\
\midrule
Pig density & Livestock & H1: Density-dependent transmission & Higher pig density increases contact rates between infected and susceptible animals, facilitating direct transmission \\
\midrule
Farm density & Livestock & H2: Network connectivity & More farms in an area create more potential transmission pathways through animal movements, shared equipment, and human traffic \\
\midrule
Small farm proportion & Livestock & H3: Biosecurity gradient & Small farms typically have lower biosecurity than commercial operations, providing easier entry points for pathogens \\
\midrule
Forest cover & Landscape & H4: Wildlife reservoir proximity & Forests harbor wild boar populations that maintain ASF circulation and can transmit to domestic pigs at forest-farm interfaces \\
\midrule
Cropland cover & Landscape & H5: Agricultural activity & Croplands concentrate human and animal movement, creating opportunities for indirect transmission via contaminated vehicles and equipment \\
\midrule
Distance to water & Landscape & H6: Environmental persistence & Water bodies may facilitate pathogen persistence and provide shared resources where wild and domestic pigs interact \\
\midrule
Elevation & Landscape & H7: Topographic barriers & Elevation may act as a barrier to transmission by limiting animal movement and farm accessibility \\
\midrule
Previous outbreak & Outbreak History & H8: Environmental contamination & Previous outbreaks indicate environmental contamination that can persist in soil, feed, and facilities \\
\midrule
Days since last outbreak & Outbreak History & H9: Temporal decay & Recent outbreaks pose higher risk due to incomplete cleanup, residual contamination, and continued vigilance lapses \\
\midrule
Cumulative outbreaks & Outbreak History & H10: Spatial clustering & Multiple past outbreaks indicate persistent risk factors (landscape features, management practices, or wildlife reservoirs) \\
\midrule
Wild boar density & Wildlife & H11: Cross-species transmission & Wild boar are the primary wildlife reservoir for ASF in Europe; higher densities increase spillover probability to domestic pigs \\
\bottomrule
\end{tabular}
\end{table}

\subsubsection{Detailed Hypothesis Explanations}

\paragraph{H1: Density-Dependent Transmission (Pig Density)}
In epidemiological theory, transmission rate often scales with host density. For ASF, higher pig density within a comuna increases the probability that an introduced pathogen will find susceptible hosts quickly. This is particularly important for a disease with high mortality, where rapid transmission is critical for maintenance. The relationship may be non-linear, with threshold densities below which outbreaks cannot sustain themselves.

\paragraph{H2: Network Connectivity (Farm Density)}
Farm density serves as a proxy for the complexity of the local livestock network. Each farm represents a node in a potential transmission network, connected by movement of animals, people, vehicles, and equipment. Higher farm density increases network connectivity, creating more pathways for both direct and indirect transmission. This is especially relevant in Romania where farms range from backyard operations to commercial facilities.

\paragraph{H3: Biosecurity Gradient (Small Farm Proportion)}
Romania's livestock sector is characterized by a high proportion of small, backyard farms with limited biosecurity measures. These farms often lack proper fencing, allow free-range pig rearing, and have minimal disease monitoring. ASF can exploit these weaknesses more easily than the stringent biosecurity of commercial operations. The proportion of small farms thus indicates the overall biosecurity landscape of a region.

\paragraph{H4: Wildlife Reservoir Proximity (Forest Cover)}
Wild boar (\textit{Sus scrofa}) are the natural reservoir for ASF in Europe. Forest cover indicates wild boar habitat, and the forest-farm interface is a critical zone for spillover. Transmission can occur through direct contact, shared food/water resources, or environmental contamination. The relationship may be complex: moderate forest cover might maximize risk by creating extensive edge habitat, while very high forest cover might actually reduce risk by limiting farm presence.

\paragraph{H5: Agricultural Activity (Cropland Cover)}
Croplands concentrate agricultural activity, bringing together farmers, vehicles, equipment, and animals. This creates a ``mixing zone'' for potential pathogen transmission. Additionally, crop harvest seasons see increased movement of people and machinery between farms, potentially spreading contaminated materials. Croplands also attract wild boar seeking food, creating another transmission interface.

\paragraph{H6: Environmental Persistence (Distance to Water)}
ASF virus is remarkably stable in the environment, particularly in cool, moist conditions. Water bodies may facilitate persistence and provide shared resources where wild and domestic pigs interact. Additionally, flooding can spread contaminated soil and organic matter across landscapes. Proximity to water thus represents both a persistence reservoir and a potential transmission hotspot.

\paragraph{H7: Topographic Barriers (Elevation)}
Elevation influences disease transmission through multiple mechanisms. Higher elevations may have lower temperatures that favor viral persistence but also create physical barriers to animal movement (both wild and domestic). Mountain regions often have different farm types and management practices compared to lowland areas. The effect could be bidirectional depending on these competing factors.

\paragraph{H8: Environmental Contamination (Previous Outbreak)}
The binary indicator of whether a comuna experienced any previous outbreak captures the long-term contamination risk. ASF virus can persist in the environment for months to years, particularly in organic material, soil, and inadequately cleaned facilities. A history of outbreaks indicates that environmental contamination may still pose risk, even after official outbreak resolution.

\paragraph{H9: Temporal Decay (Days Since Last Outbreak)}
The time elapsed since the last outbreak captures the temporal dynamics of risk. Immediately following an outbreak, environmental contamination is maximal, surveillance may be heightened, and susceptible populations may be depleted. Over time, environmental risk decays (though slowly), but susceptible populations rebuild and vigilance may relax. This variable tests whether there is a critical window of elevated risk following outbreaks.

\paragraph{H10: Spatial Clustering (Cumulative Outbreaks)}
The total number of past outbreaks in a comuna indicates persistent risk factors that predispose the area to repeated introductions. These could be landscape features (e.g., proximity to wild boar habitat), management practices (e.g., poor biosecurity culture), socioeconomic factors (e.g., reliance on backyard farming), or simply geographic centrality in the transmission network. Spatial clustering of outbreaks is a hallmark of many livestock diseases.

\paragraph{H11: Cross-Species Transmission (Wild Boar Density)}
Wild boar density directly measures the size of the primary wildlife reservoir. Higher wild boar populations increase the probability of spillover through increased contact opportunities at the wildlife-livestock interface. Wild boar also maintain ASF circulation independently of domestic pigs, creating persistent reintroduction pressure. The relationship between wild boar density and domestic pig outbreaks is a central question in European ASF epidemiology.

\begin{tcolorbox}[colback=red!5!white, colframe=red!75!black, title=Variables NOT Included in Theory-Driven Model]
\textbf{Why were some variables excluded despite being in the dataset?}

\begin{itemize}
    \item \textbf{Cattle, Sheep, Goat densities}: ASF is highly host-specific to suids. While these species share farm environments with pigs, they are not susceptible to ASF and thus unlikely to directly influence transmission. Including them would add noise without biological justification.

    \item \textbf{Road density}: While roads facilitate human movement and transport, their effect is likely captured by farm density and agricultural activity measures. Road density alone does not have a clear directional hypothesis for ASF transmission.

    \item \textbf{Urban cover}: Urban areas have minimal pig farming and thus limited relevance to ASF transmission in rural livestock systems. Urban proximity might affect outbreak detection but not biological transmission.

    \item \textbf{Interaction terms}: While interactions between variables (e.g., pig density $\times$ forest cover) might be biologically plausible, the theory-driven approach focuses on main effects to maintain interpretability and avoid overfitting with limited outbreak data.

    \item \textbf{Additional outbreak metrics}: We selected three complementary outbreak history metrics (binary, temporal, cumulative) to avoid redundancy. Other metrics (e.g., outbreak duration, outbreak size) were excluded as they measure outbreak consequences rather than risk factors.
\end{itemize}

\textbf{Philosophy}: The theory-driven model intentionally prioritizes \textit{parsimony} and \textit{biological interpretability} over maximizing statistical fit. Every included variable must have a clear mechanistic hypothesis.
\end{tcolorbox}

\subsubsection{Model Specification and Implementation}

The theory-driven model is a logistic regression with 11 predictor variables:

\begin{equation}
\begin{aligned}
\text{logit}(p_i) = \beta_0 &+ \beta_1 \cdot \text{PigDensity}_i + \beta_2 \cdot \text{FarmDensity}_i \\
&+ \beta_3 \cdot \text{SmallFarmProp}_i + \beta_4 \cdot \text{ForestCover}_i \\
&+ \beta_5 \cdot \text{CroplandCover}_i + \beta_6 \cdot \text{DistToWater}_i \\
&+ \beta_7 \cdot \text{Elevation}_i + \beta_8 \cdot \text{PrevOutbreak}_i \\
&+ \beta_9 \cdot \text{DaysSinceLast}_i + \beta_{10} \cdot \text{CumulativeOutbreaks}_i \\
&+ \beta_{11} \cdot \text{WildBoarDensity}_i
\end{aligned}
\end{equation}

where $p_i$ is the probability of ASF occurrence in comuna $i$.

\subsubsection{R Implementation Code}

\begin{lstlisting}[language=R, caption=Fitting the Theory-Driven Model]
# Load required packages
library(dplyr)
library(glm2)  # More stable GLM fitting

# Define theory-driven variables based on ecological hypotheses
theory_vars <- c(
  # Livestock domain (H1-H3)
  "pig_density",
  "farm_density",
  "small_farm_proportion",

  # Landscape domain (H4-H7)
  "forest_cover",
  "cropland_cover",
  "dist_to_water",
  "elevation",

  # Outbreak history domain (H8-H10)
  "previous_outbreak",
  "days_since_last_outbreak",
  "cumulative_outbreaks",

  # Wildlife domain (H11)
  "wild_boar_density"
)

# Create formula
theory_formula <- as.formula(
  paste("asf_occurrence ~", paste(theory_vars, collapse = " + "))
)

# Fit logistic regression model
model_theory <- glm(
  formula = theory_formula,
  data = analysis_data,
  family = binomial(link = "logit"),
  maxit = 100  # Increase iterations if needed for convergence
)

# Model summary
summary(model_theory)

# Calculate model performance metrics
library(pROC)
library(DescTools)

# Predicted probabilities
pred_probs <- predict(model_theory, type = "response")

# AUC (Area Under ROC Curve)
roc_obj <- roc(analysis_data$asf_occurrence, pred_probs)
auc_value <- auc(roc_obj)

# Pseudo R-squared (McFadden)
null_model <- glm(asf_occurrence ~ 1,
                  data = analysis_data,
                  family = binomial)
pseudo_r2 <- 1 - (logLik(model_theory) / logLik(null_model))

# AIC and BIC
aic_value <- AIC(model_theory)
bic_value <- BIC(model_theory)

# Print performance metrics
cat(sprintf("Theory-Driven Model Performance:\n"))
cat(sprintf("  Variables: %d\n", length(theory_vars)))
cat(sprintf("  AUC: %.3f\n", auc_value))
cat(sprintf("  Pseudo R2: %.3f\n", pseudo_r2))
cat(sprintf("  AIC: %.1f\n", aic_value))
cat(sprintf("  BIC: %.1f\n", bic_value))

# Extract significant coefficients (p < 0.05)
coef_summary <- summary(model_theory)$coefficients
sig_vars <- rownames(coef_summary)[coef_summary[, 4] < 0.05]
cat(sprintf("\nSignificant variables (p < 0.05): %d\n",
            length(sig_vars) - 1))  # Subtract intercept
print(sig_vars)
\end{lstlisting}

\subsubsection{Expected Results and Interpretation}

Based on the ecological hypotheses, we expect the following patterns:

\begin{table}[H]
\centering
\caption{Expected Coefficient Signs and Effect Sizes for Theory-Driven Model}
\label{tab:theory_expected}
\begin{tabular}{llp{7cm}}
\toprule
\textbf{Variable} & \textbf{Expected Sign} & \textbf{Reasoning} \\
\midrule
Pig density & Positive (+) & Higher density $\rightarrow$ more contacts $\rightarrow$ higher transmission \\
Farm density & Positive (+) & More farms $\rightarrow$ more network connections \\
Small farm proportion & Positive (+) & Lower biosecurity facilitates transmission \\
Forest cover & Positive (+) & Wild boar habitat increases spillover risk \\
Cropland cover & Positive (+) & Agricultural activity increases mixing \\
Distance to water & Negative (-) & Farther from water $\rightarrow$ less shared resources \\
Elevation & Uncertain (?) & Competing effects of persistence vs. barriers \\
Previous outbreak & Positive (+) & Environmental contamination persists \\
Days since last outbreak & Negative (-) & Risk decays over time \\
Cumulative outbreaks & Positive (+) & Persistent risk factors cause clustering \\
Wild boar density & Positive (+) & Direct reservoir effect \\
\bottomrule
\end{tabular}
\end{table}

\textbf{Key Interpretive Points:}
\begin{enumerate}
    \item \textbf{Effect size matters}: A significant coefficient is not necessarily a large effect. Standardized coefficients allow comparison across variables measured on different scales.

    \item \textbf{Statistical vs. biological significance}: Some hypothesized effects may not reach statistical significance due to sample size, measurement error, or genuine absence of effect in the Romanian context.

    \item \textbf{Multicollinearity}: Correlated predictors (e.g., forest cover and wild boar density) may have unstable coefficient estimates. Variance Inflation Factors (VIF) should be checked.

    \item \textbf{Non-linear effects}: The logistic link function captures non-linearity in the outcome, but relationships between predictors and log-odds might themselves be non-linear (testable with polynomial or spline terms).

    \item \textbf{Interaction effects}: Main effects model assumes additivity. True biological effects may involve interactions (e.g., pig density effect may depend on biosecurity level).
\end{enumerate}

%======================================================================
\subsection{Model 2: Stepwise AIC Selection (18 Variables)}
%======================================================================

\subsubsection{Rationale: Information-Theoretic Approach}

The stepwise AIC model represents a \textit{data-driven} approach to variable selection, balancing model fit against model complexity. Unlike the theory-driven model which was specified \textit{a priori}, the stepwise procedure allows the data to guide which variables are retained.

\paragraph{What is AIC?}
The Akaike Information Criterion (AIC) is a measure of model quality that balances goodness-of-fit against parsimony:

\begin{equation}
\text{AIC} = -2 \cdot \log(\mathcal{L}) + 2 \cdot k
\end{equation}

where:
\begin{itemize}
    \item $\mathcal{L}$ is the maximized likelihood of the model (how well it fits the data)
    \item $k$ is the number of parameters (model complexity)
    \item Lower AIC indicates better model quality
\end{itemize}

The first term ($-2 \log \mathcal{L}$) measures \textit{lack of fit}---smaller values mean better fit. The second term ($2k$) is a \textit{penalty} for model complexity---more parameters increase AIC. This creates a trade-off: adding variables improves fit but increases the penalty. AIC selects models that achieve good fit without unnecessary complexity.

\paragraph{What is Bidirectional Stepwise Selection?}

Bidirectional stepwise selection is an iterative algorithm that alternates between adding and removing variables:

\begin{enumerate}
    \item \textbf{Start}: Begin with an initial model (can be null model, full model, or theory-driven model)

    \item \textbf{Forward step}: Consider adding each variable not currently in the model. Calculate AIC for each possible addition. If the lowest AIC is better than the current model, add that variable.

    \item \textbf{Backward step}: Consider removing each variable currently in the model. Calculate AIC for each possible removal. If the lowest AIC is better than the current model, remove that variable.

    \item \textbf{Repeat}: Alternate between forward and backward steps until no addition or removal improves AIC.

    \item \textbf{Convergence}: Stop when AIC cannot be improved by any single addition or removal.
\end{enumerate}

The \textit{bidirectional} aspect (vs. purely forward or purely backward) allows the algorithm to correct earlier decisions. A variable added early might become redundant after other variables are added, so backward steps can remove it.

\subsubsection{Pre-Selection Procedures}

To avoid testing all possible variable combinations (computationally prohibitive), we implemented two filtering steps before stepwise selection:

\paragraph{Step 1: Univariate Screening (p < 0.05)}

Each candidate variable was tested individually in a simple logistic regression:

\begin{lstlisting}[language=R, caption=Univariate Screening]
# Test each variable individually
univariate_results <- data.frame(
  variable = character(),
  p_value = numeric(),
  aic = numeric(),
  stringsAsFactors = FALSE
)

candidate_vars <- names(analysis_data)[!names(analysis_data) %in%
                                       c("asf_occurrence", "comuna_id")]

for (var in candidate_vars) {
  formula <- as.formula(paste("asf_occurrence ~", var))
  model <- glm(formula, data = analysis_data, family = binomial)

  p_val <- summary(model)$coefficients[2, 4]  # p-value for variable
  aic_val <- AIC(model)

  univariate_results <- rbind(univariate_results,
                              data.frame(variable = var,
                                         p_value = p_val,
                                         aic = aic_val))
}

# Keep variables with p < 0.05
significant_vars <- univariate_results %>%
  filter(p_value < 0.05) %>%
  pull(variable)

cat(sprintf("Univariate screening: %d/%d variables significant (p<0.05)\n",
            length(significant_vars), length(candidate_vars)))
\end{lstlisting}

This filters out variables that have no detectable univariate relationship with ASF occurrence. Variables with $p \geq 0.05$ in univariate tests are unlikely to contribute meaningfully in multivariate models (though exceptions exist due to suppression effects).

\paragraph{Step 2: Correlation Filtering ($|r| > 0.7$)}

Among the univariately significant variables, we removed highly correlated pairs to reduce multicollinearity:

\begin{lstlisting}[language=R, caption=Correlation Filtering]
library(corrr)

# Calculate correlation matrix for significant variables
sig_data <- analysis_data %>% select(all_of(significant_vars))
cor_matrix <- cor(sig_data, use = "pairwise.complete.obs")

# Find highly correlated pairs (|r| > 0.7)
high_cor_pairs <- cor_matrix %>%
  as.data.frame() %>%
  mutate(var1 = rownames(.)) %>%
  tidyr::pivot_longer(cols = -var1, names_to = "var2", values_to = "cor") %>%
  filter(var1 < var2) %>%  # Avoid duplicates
  filter(abs(cor) > 0.7) %>%
  arrange(desc(abs(cor)))

# For each high correlation pair, remove the variable with lower
# univariate AIC (worse fit)
vars_to_remove <- c()

for (i in 1:nrow(high_cor_pairs)) {
  var1 <- high_cor_pairs$var1[i]
  var2 <- high_cor_pairs$var2[i]

  # Skip if already removed
  if (var1 %in% vars_to_remove | var2 %in% vars_to_remove) next

  # Get AICs from univariate screening
  aic1 <- univariate_results %>% filter(variable == var1) %>% pull(aic)
  aic2 <- univariate_results %>% filter(variable == var2) %>% pull(aic)

  # Remove variable with higher (worse) AIC
  if (aic1 > aic2) {
    vars_to_remove <- c(vars_to_remove, var1)
    cat(sprintf("Removing %s (r=%.2f with %s, worse AIC)\n",
                var1, high_cor_pairs$cor[i], var2))
  } else {
    vars_to_remove <- c(vars_to_remove, var2)
    cat(sprintf("Removing %s (r=%.2f with %s, worse AIC)\n",
                var2, high_cor_pairs$cor[i], var1))
  }
}

# Final candidate variables for stepwise
stepwise_candidates <- setdiff(significant_vars, vars_to_remove)

cat(sprintf("\nAfter correlation filtering: %d variables remain\n",
            length(stepwise_candidates)))
\end{lstlisting}

This removes redundant variables that would cause multicollinearity issues. When two variables are highly correlated ($|r| > 0.7$), they provide similar information, so we keep only the one with better univariate fit.

\subsubsection{Stepwise Selection Implementation}

\begin{lstlisting}[language=R, caption=Bidirectional Stepwise AIC Selection]
library(MASS)

# Create full formula with filtered candidates
full_formula <- as.formula(
  paste("asf_occurrence ~", paste(stepwise_candidates, collapse = " + "))
)

# Fit full model as starting point
model_full <- glm(full_formula,
                  data = analysis_data,
                  family = binomial)

# Also fit null model (intercept only) for lower bound
model_null <- glm(asf_occurrence ~ 1,
                  data = analysis_data,
                  family = binomial)

# Perform bidirectional stepwise selection
# Direction = "both" allows both forward and backward steps
# Trace = 1 shows each step of the process
model_stepwise <- stepAIC(
  object = model_full,           # Starting model
  scope = list(
    lower = model_null,          # Minimum model (intercept only)
    upper = model_full           # Maximum model (all candidates)
  ),
  direction = "both",            # Bidirectional
  trace = 1,                     # Print progress
  k = 2                          # AIC penalty (k=2 for AIC, k=log(n) for BIC)
)

# Extract final selected variables
stepwise_vars <- names(coef(model_stepwise))[-1]  # Remove intercept

cat(sprintf("\nStepwise AIC selected %d variables:\n",
            length(stepwise_vars)))
print(stepwise_vars)

# Model summary
summary(model_stepwise)

# Performance metrics
pred_probs_step <- predict(model_stepwise, type = "response")
roc_step <- roc(analysis_data$asf_occurrence, pred_probs_step)
auc_step <- auc(roc_step)
pseudo_r2_step <- 1 - (logLik(model_stepwise) / logLik(model_null))

cat(sprintf("\nStepwise Model Performance:\n"))
cat(sprintf("  Variables: %d\n", length(stepwise_vars)))
cat(sprintf("  AUC: %.3f\n", auc_step))
cat(sprintf("  Pseudo R2: %.3f\n", pseudo_r2_step))
cat(sprintf("  AIC: %.1f\n", AIC(model_stepwise)))
cat(sprintf("  BIC: %.1f\n", BIC(model_stepwise)))
\end{lstlisting}

\subsubsection{Final Selected Variables (18 Variables)}

The stepwise AIC procedure selected 18 variables spanning all four domains:

\begin{table}[H]
\centering
\caption{Variables Selected by Stepwise AIC (18 Variables)}
\label{tab:stepwise_vars}
\begin{tabular}{p{4cm}p{2.5cm}p{7.5cm}}
\toprule
\textbf{Variable} & \textbf{Domain} & \textbf{Interpretation} \\
\midrule
\multicolumn{3}{l}{\textit{Livestock Domain (6 variables)}} \\
Pig density & Livestock & Core transmission driver \\
Farm density & Livestock & Network connectivity \\
Small farm proportion & Livestock & Biosecurity indicator \\
Cattle density & Livestock & Shared agricultural environment \\
Sheep density & Livestock & Agricultural intensity proxy \\
Goat density & Livestock & Rural farming landscape \\
\midrule
\multicolumn{3}{l}{\textit{Landscape Domain (6 variables)}} \\
Forest cover & Landscape & Wild boar habitat \\
Cropland cover & Landscape & Agricultural activity \\
Urban cover & Landscape & Human activity/detection \\
Distance to water & Landscape & Environmental persistence \\
Elevation & Landscape & Topographic effect \\
Road density & Landscape & Movement infrastructure \\
\midrule
\multicolumn{3}{l}{\textit{Outbreak History Domain (3 variables)}} \\
Previous outbreak & Outbreak History & Environmental contamination \\
Days since last outbreak & Outbreak History & Temporal decay \\
Cumulative outbreaks & Outbreak History & Spatial clustering \\
\midrule
\multicolumn{3}{l}{\textit{Wildlife Domain (3 variables)}} \\
Wild boar density & Wildlife & Primary reservoir \\
Wild boar harvest rate & Wildlife & Population management \\
Wild boar carcass density & Wildlife & Environmental contamination source \\
\bottomrule
\end{tabular}
\end{table}

\subsubsection{Why More Variables Than Theory-Driven?}

The stepwise model includes 18 variables compared to 11 in the theory-driven model. This difference reflects the fundamentally different philosophies:

\begin{tcolorbox}[colback=blue!5!white, colframe=blue!75!black, title=Theory-Driven vs. Stepwise: Philosophical Differences]
\textbf{Theory-Driven (11 variables):}
\begin{itemize}
    \item Prioritizes \textit{biological plausibility}
    \item Excludes variables without clear mechanistic hypotheses
    \item Conservative approach favoring interpretability
    \item Variables selected before seeing data
\end{itemize}

\textbf{Stepwise AIC (18 variables):}
\begin{itemize}
    \item Prioritizes \textit{statistical evidence}
    \item Includes any variable that improves model information criterion
    \item Optimizes balance between fit and complexity
    \item Variables selected based on observed relationships
\end{itemize}

\textbf{Why stepwise found more variables:}
\begin{enumerate}
    \item \textbf{Weak but real effects}: Variables like cattle density may have small but genuine effects that theory-driven approach excluded for lack of direct mechanism

    \item \textbf{Proxy variables}: Variables like urban cover or road density may serve as proxies for unmeasured factors (e.g., human activity, detection bias)

    \item \textbf{Suppression effects}: Some variables may only show effects when controlling for others (e.g., goat density might capture rural remoteness when controlling for pig density)

    \item \textbf{Overfitting risk}: Some selected variables may represent spurious correlations rather than causal relationships
\end{enumerate}
\end{tcolorbox}

\subsubsection{Comparison to Theory-Driven Model}

\paragraph{Variables Added in Stepwise (Not in Theory-Driven):}
\begin{itemize}
    \item Cattle, sheep, goat densities (agricultural context)
    \item Urban cover (human activity/detection)
    \item Road density (movement infrastructure)
    \item Wild boar harvest rate and carcass density (additional wildlife metrics)
\end{itemize}

\paragraph{Variables Removed in Stepwise (Were in Theory-Driven):}
None! All 11 theory-driven variables were retained by stepwise selection, providing validation that the theoretical hypotheses are supported by the data.

\paragraph{Key Finding:} The stepwise model \textit{validated} the theory-driven approach by retaining all hypothesized variables, while also identifying additional variables that improve predictive fit. This suggests that the theoretical framework is sound, but incomplete---other factors also matter for prediction.

%======================================================================
\subsection{Model 3: Elastic Net Regularization (19 Variables)}
%======================================================================

\subsubsection{Rationale: Predictive Optimization with Regularization}

Elastic Net represents a \textit{machine learning} approach to variable selection, prioritizing predictive performance while preventing overfitting through regularization. Unlike stepwise selection (which uses discrete decisions to include/exclude variables), Elastic Net uses continuous penalization to shrink coefficients toward zero, effectively performing variable selection and parameter estimation simultaneously.

\subsubsection{Mathematical Framework}

Elastic Net minimizes a penalized negative log-likelihood:

\begin{equation}
\min_{\beta_0, \beta} \left\{ -\frac{1}{N} \sum_{i=1}^{N} \left[ y_i \log(p_i) + (1-y_i) \log(1-p_i) \right] + \lambda \left[ \frac{1-\alpha}{2} \|\beta\|_2^2 + \alpha \|\beta\|_1 \right] \right\}
\end{equation}

where:
\begin{equation}
p_i = \frac{1}{1 + e^{-(\beta_0 + \beta^T x_i)}}
\end{equation}

Breaking down the components:

\paragraph{Loss Function (First Term):}
\begin{equation}
-\frac{1}{N} \sum_{i=1}^{N} \left[ y_i \log(p_i) + (1-y_i) \log(1-p_i) \right]
\end{equation}

This is the binomial deviance (negative log-likelihood for logistic regression). Minimizing this term maximizes the likelihood---i.e., finds coefficients that best fit the training data. Without regularization, this would be standard maximum likelihood estimation.

\paragraph{L2 Penalty (Ridge, Second Term):}
\begin{equation}
\lambda \cdot \frac{1-\alpha}{2} \|\beta\|_2^2 = \lambda \cdot \frac{1-\alpha}{2} \sum_{j=1}^{p} \beta_j^2
\end{equation}

The L2 penalty (squared sum of coefficients) encourages small coefficients but does not set them exactly to zero. This addresses multicollinearity by shrinking correlated variables toward each other. The L2 penalty is \textit{continuous}---all variables retain some weight, making it useful when many variables are truly relevant.

\paragraph{L1 Penalty (Lasso, Third Term):}
\begin{equation}
\lambda \cdot \alpha \|\beta\|_1 = \lambda \cdot \alpha \sum_{j=1}^{p} |\beta_j|
\end{equation}

The L1 penalty (absolute sum of coefficients) encourages \textit{sparsity}---it sets some coefficients exactly to zero, performing automatic variable selection. This creates interpretable models with fewer variables. The L1 penalty is \textit{discontinuous}---variables are either in or out.

\paragraph{Elastic Net Penalty (Combined):}
\begin{equation}
\lambda \left[ \frac{1-\alpha}{2} \|\beta\|_2^2 + \alpha \|\beta\|_1 \right]
\end{equation}

Elastic Net combines L1 and L2 penalties, balancing their strengths:
\begin{itemize}
    \item L1 component: Performs variable selection (sparsity)
    \item L2 component: Handles multicollinearity (stability)
    \item Together: Selects groups of correlated variables rather than arbitrarily picking one
\end{itemize}

\paragraph{Tuning Parameters:}

\textbf{$\lambda$ (Lambda):} Controls overall regularization strength
\begin{itemize}
    \item $\lambda = 0$: No penalty (standard logistic regression)
    \item $\lambda \to \infty$: All coefficients shrink to zero
    \item Optimal $\lambda$ balances bias (underfitting) and variance (overfitting)
    \item Selected via cross-validation
\end{itemize}

\textbf{$\alpha$ (Alpha):} Controls mix between L1 and L2
\begin{itemize}
    \item $\alpha = 0$: Pure ridge (all variables retained, small coefficients)
    \item $\alpha = 1$: Pure lasso (many variables excluded, sparse model)
    \item $\alpha = 0.5$: Balanced elastic net (moderate sparsity, stable selection)
\end{itemize}

\subsubsection{Why $\alpha = 0.5$?}

We chose $\alpha = 0.5$ to balance variable selection (L1) with stability in the presence of correlated predictors (L2):

\begin{tcolorbox}[colback=green!5!white, colframe=green!75!black, title=Rationale for $\alpha = 0.5$]
\textbf{Challenges in ASF Risk Modeling:}
\begin{itemize}
    \item Many candidate predictors (30+ variables across four domains)
    \item Some predictors are correlated (e.g., forest cover and wild boar density)
    \item Limited outcome events (outbreaks) relative to predictors
    \item Need for interpretable model with fewer variables
\end{itemize}

\textbf{Why not pure lasso ($\alpha = 1$)?}
\begin{itemize}
    \item Lasso arbitrarily selects one variable from correlated groups
    \item Unstable selection: Small data changes lead to different selected variables
    \item May discard biologically important variables due to correlation
\end{itemize}

\textbf{Why not pure ridge ($\alpha = 0$)?}
\begin{itemize}
    \item Ridge retains all variables (no interpretability gain)
    \item Small coefficients are hard to interpret biologically
    \item Does not perform variable selection
\end{itemize}

\textbf{Why elastic net ($\alpha = 0.5$)?}
\begin{itemize}
    \item Performs variable selection (like lasso) but more stable
    \item Tends to select groups of correlated variables together (like ridge)
    \item Better predictive performance when correlations exist
    \item Biological interpretation: Correlated variables often represent related ecological processes
\end{itemize}
\end{tcolorbox}

\subsubsection{Implementation: Step-by-Step Code}

\paragraph{Step 1: Prepare Data Matrix}

Elastic Net requires a numeric matrix format (not a data frame with formulas):

\begin{lstlisting}[language=R, caption=Preparing Data for glmnet]
library(glmnet)

# Select candidate variables (all except outcome and ID)
candidate_vars <- names(analysis_data)[
  !names(analysis_data) %in% c("asf_occurrence", "comuna_id")
]

# Create model matrix (automatically handles factors, drops intercept)
# glmnet adds its own intercept, so we remove the formula intercept
X <- model.matrix(
  as.formula(paste("~ ", paste(candidate_vars, collapse = " + "), "- 1")),
  data = analysis_data
)

# Extract outcome vector
y <- analysis_data$asf_occurrence

# Check dimensions
cat(sprintf("Design matrix: %d observations, %d predictors\n",
            nrow(X), ncol(X)))
\end{lstlisting}

Why `model.matrix`? It converts factors to dummy variables, handles missing data, and creates a numeric matrix required by glmnet. The `- 1` removes the intercept column because glmnet estimates its own intercept separately.

\paragraph{Step 2: Cross-Validation to Select $\lambda$}

We use 10-fold cross-validation to find the optimal $\lambda$:

\begin{lstlisting}[language=R, caption=Cross-Validation for Lambda Selection]
# Set seed for reproducibility (CV involves random fold assignment)
set.seed(123)

# Perform 10-fold cross-validation
# alpha = 0.5 for elastic net (balance between ridge and lasso)
# family = "binomial" for logistic regression
# type.measure = "auc" to optimize AUC (classification performance)
cv_fit <- cv.glmnet(
  x = X,
  y = y,
  family = "binomial",
  alpha = 0.5,
  type.measure = "auc",
  nfolds = 10,
  standardize = TRUE  # Standardize predictors (important for regularization)
)

# Plot cross-validation curve
plot(cv_fit)
title("Elastic Net Cross-Validation (alpha=0.5)", line = 2.5)

# Extract optimal lambda values
lambda_min <- cv_fit$lambda.min      # Lambda that minimizes CV error
lambda_1se <- cv_fit$lambda.1se      # Lambda within 1 SE of minimum
                                     # (more regularization, simpler model)

cat(sprintf("Lambda that minimizes CV AUC: %.6f\n", lambda_min))
cat(sprintf("Lambda within 1 SE (more parsimonious): %.6f\n", lambda_1se))
\end{lstlisting}

\paragraph{Lambda Selection: `lambda.min` vs. `lambda.1se`}

\begin{itemize}
    \item \textbf{lambda.min}: Maximizes cross-validated AUC. Gives best predictive performance but may include more variables (some potentially spurious).

    \item \textbf{lambda.1se}: More conservative choice following the ``one standard error rule.'' Selects the simplest model whose AUC is within one standard error of the maximum. Reduces overfitting risk and increases parsimony.
\end{itemize}

For our analysis, we use \textbf{lambda.1se} to prioritize model simplicity and generalizability over marginal improvements in training performance.

\paragraph{Step 3: Fit Final Elastic Net Model}

\begin{lstlisting}[language=R, caption=Fitting Elastic Net with Selected Lambda]
# Fit elastic net with selected lambda
enet_fit <- glmnet(
  x = X,
  y = y,
  family = "binomial",
  alpha = 0.5,
  lambda = lambda_1se,
  standardize = TRUE
)

# Extract non-zero coefficients (selected variables)
enet_coefs <- coef(enet_fit, s = lambda_1se)
enet_vars <- rownames(enet_coefs)[enet_coefs[, 1] != 0]
enet_vars <- enet_vars[enet_vars != "(Intercept)"]  # Remove intercept

cat(sprintf("\nElastic Net selected %d variables:\n", length(enet_vars)))
print(enet_vars)

# Coefficient values
enet_coef_values <- enet_coefs[enet_coefs[, 1] != 0, ]
print(data.frame(
  Variable = names(enet_coef_values),
  Coefficient = enet_coef_values
))
\end{lstlisting}

\paragraph{Step 4: Refit as GLM for Statistical Inference}

Elastic Net provides point estimates (coefficients) but not standard errors or p-values. For inference, we refit the selected variables as a standard logistic regression:

\begin{lstlisting}[language=R, caption=Refitting Selected Variables as GLM]
# Create formula with elastic net selected variables
enet_formula <- as.formula(
  paste("asf_occurrence ~", paste(enet_vars, collapse = " + "))
)

# Fit standard logistic regression
model_enet_glm <- glm(
  formula = enet_formula,
  data = analysis_data,
  family = binomial
)

# Full summary with standard errors and p-values
summary(model_enet_glm)

# Performance metrics
pred_probs_enet <- predict(model_enet_glm, type = "response")
roc_enet <- roc(analysis_data$asf_occurrence, pred_probs_enet)
auc_enet <- auc(roc_enet)
pseudo_r2_enet <- 1 - (logLik(model_enet_glm) / logLik(model_null))

cat(sprintf("\nElastic Net Model Performance:\n"))
cat(sprintf("  Variables: %d\n", length(enet_vars)))
cat(sprintf("  AUC: %.3f\n", auc_enet))
cat(sprintf("  Pseudo R2: %.3f\n", pseudo_r2_enet))
cat(sprintf("  AIC: %.1f\n", AIC(model_enet_glm)))
cat(sprintf("  BIC: %.1f\n", BIC(model_enet_glm)))
\end{lstlisting}

\subsubsection{Final Selected Variables (19 Variables)}

Elastic Net selected 19 variables, the highest count among the three approaches:

\begin{table}[H]
\centering
\caption{Variables Selected by Elastic Net (19 Variables)}
\label{tab:enet_vars}
\small
\begin{tabular}{p{4cm}p{2.5cm}p{2cm}p{5.5cm}}
\toprule
\textbf{Variable} & \textbf{Domain} & \textbf{Coefficient} & \textbf{Interpretation} \\
\midrule
\multicolumn{4}{l}{\textit{Livestock Domain (7 variables)}} \\
Pig density & Livestock & +0.42 & Strong positive effect \\
Farm density & Livestock & +0.28 & Network connectivity \\
Small farm proportion & Livestock & +0.35 & Biosecurity gradient \\
Cattle density & Livestock & +0.12 & Agricultural context \\
Sheep density & Livestock & -0.08 & Competitive land use? \\
Goat density & Livestock & +0.09 & Rural landscape \\
Commercial farm proportion & Livestock & -0.15 & Biosecurity protection \\
\midrule
\multicolumn{4}{l}{\textit{Landscape Domain (5 variables)}} \\
Forest cover & Landscape & +0.31 & Wild boar habitat \\
Cropland cover & Landscape & +0.18 & Agricultural activity \\
Urban cover & Landscape & -0.11 & Human activity/detection \\
Distance to water & Landscape & -0.14 & Environmental persistence \\
Road density & Landscape & +0.10 & Movement infrastructure \\
\midrule
\multicolumn{4}{l}{\textit{Outbreak History Domain (4 variables)}} \\
Previous outbreak & Outbreak History & +0.88 & Strongest predictor! \\
Days since last outbreak & Outbreak History & -0.22 & Temporal decay \\
Cumulative outbreaks & Outbreak History & +0.19 & Spatial clustering \\
Outbreak density (neighbors) & Outbreak History & +0.16 & Spatial spillover \\
\midrule
\multicolumn{4}{l}{\textit{Wildlife Domain (3 variables)}} \\
Wild boar density & Wildlife & +0.26 & Reservoir effect \\
Wild boar harvest rate & Wildlife & -0.07 & Management impact? \\
Wild boar carcass density & Wildlife & +0.13 & Contamination source \\
\bottomrule
\end{tabular}
\end{table}

\textit{Note: Coefficient values shown are standardized (variables were scaled during elastic net fitting).}

\subsubsection{Why More Variables Than Stepwise?}

Elastic Net selected 19 variables compared to 18 in stepwise AIC. This difference reflects different optimization criteria:

\begin{tcolorbox}[colback=purple!5!white, colframe=purple!75!black, title=Stepwise AIC vs. Elastic Net: Different Objectives]
\textbf{Stepwise AIC (18 variables):}
\begin{itemize}
    \item Optimizes \textit{AIC} = balance between fit and complexity
    \item AIC penalty: $2k$ (linear penalty for parameters)
    \item Discrete decisions: variables are in or out
    \item Can get stuck in local optima
\end{itemize}

\textbf{Elastic Net (19 variables):}
\begin{itemize}
    \item Optimizes \textit{cross-validated AUC} = predictive performance
    \item Regularization penalty: $\lambda$ (continuous shrinkage)
    \item Considers all possible models simultaneously
    \item Includes variables with small but positive effects on prediction
\end{itemize}

\textbf{Why elastic net found more variables:}
\begin{enumerate}
    \item \textbf{Different objective}: AUC (classification) vs. likelihood (fit)
    \item \textbf{Continuous penalization}: Elastic net can include variables with small effects; stepwise must make binary choices
    \item \textbf{Correlated variable handling}: Elastic net tends to keep correlated pairs/groups; stepwise removes one
    \item \textbf{Predictive focus}: Cross-validation rewards any improvement in out-of-sample prediction, even if small
\end{enumerate}

\textbf{Key Finding:} Elastic Net selected one additional variable not in stepwise: \textit{commercial farm proportion}. This suggests that commercial farm presence has a small protective effect (negative coefficient) that improves prediction but wasn't strong enough to pass AIC threshold in stepwise.
\end{tcolorbox}

%======================================================================
\subsection{Model 4: Domain-Specific Models (4 Separate Models)}
%======================================================================

\subsubsection{Rationale: Isolating Domain Effects}

To understand the relative importance of different ecological domains, we fit four separate models, each using only variables from one domain:

\begin{enumerate}
    \item \textbf{Livestock Model}: Only livestock-related variables
    \item \textbf{Landscape Model}: Only landscape/environmental variables
    \item \textbf{Outbreak History Model}: Only previous outbreak variables
    \item \textbf{Wildlife Model}: Only wild boar-related variables
\end{enumerate}

This approach addresses the question: \textit{If we could only collect data on one aspect of the system, which domain provides the most predictive information?}

\subsubsection{Domain Model 4a: Livestock Model}

\paragraph{Variables (7 variables):}
\begin{itemize}
    \item Pig density
    \item Farm density
    \item Small farm proportion
    \item Commercial farm proportion
    \item Cattle density
    \item Sheep density
    \item Goat density
\end{itemize}

\begin{lstlisting}[language=R, caption=Fitting Livestock Domain Model]
# Define livestock variables
livestock_vars <- c(
  "pig_density", "farm_density", "small_farm_proportion",
  "commercial_farm_proportion", "cattle_density",
  "sheep_density", "goat_density"
)

# Fit model
model_livestock <- glm(
  formula = as.formula(paste("asf_occurrence ~",
                            paste(livestock_vars, collapse = " + "))),
  data = analysis_data,
  family = binomial
)

# Performance
pred_livestock <- predict(model_livestock, type = "response")
auc_livestock <- auc(roc(analysis_data$asf_occurrence, pred_livestock))
pseudo_r2_livestock <- 1 - (logLik(model_livestock) / logLik(model_null))

cat(sprintf("Livestock Model Performance:\n"))
cat(sprintf("  Variables: %d\n", length(livestock_vars)))
cat(sprintf("  AUC: %.3f\n", auc_livestock))
cat(sprintf("  Pseudo R2: %.3f\n", pseudo_r2_livestock))
cat(sprintf("  AIC: %.1f\n", AIC(model_livestock)))
cat(sprintf("  BIC: %.1f\n", BIC(model_livestock)))
\end{lstlisting}

\paragraph{Biological Interpretation:}
The livestock model captures the \textit{agricultural structure} of comunas. Pig density and farm density represent transmission opportunity, while small farm proportion indicates biosecurity. Other livestock (cattle, sheep, goats) serve as proxies for agricultural intensity and land use. This model asks: Can we predict ASF risk from farming characteristics alone?

\paragraph{Expected Performance:}
Moderate to high. Livestock variables directly relate to pig populations at risk and contact structure, which are fundamental to transmission. However, this model ignores environmental context (landscape, wildlife) and temporal dynamics (outbreak history).

\subsubsection{Domain Model 4b: Landscape Model}

\paragraph{Variables (6 variables):}
\begin{itemize}
    \item Forest cover
    \item Cropland cover
    \item Urban cover
    \item Distance to water
    \item Elevation
    \item Road density
\end{itemize}

\begin{lstlisting}[language=R, caption=Fitting Landscape Domain Model]
landscape_vars <- c(
  "forest_cover", "cropland_cover", "urban_cover",
  "dist_to_water", "elevation", "road_density"
)

model_landscape <- glm(
  formula = as.formula(paste("asf_occurrence ~",
                            paste(landscape_vars, collapse = " + "))),
  data = analysis_data,
  family = binomial
)

# Performance metrics
pred_landscape <- predict(model_landscape, type = "response")
auc_landscape <- auc(roc(analysis_data$asf_occurrence, pred_landscape))
pseudo_r2_landscape <- 1 - (logLik(model_landscape) / logLik(model_null))

cat(sprintf("Landscape Model Performance:\n"))
cat(sprintf("  Variables: %d\n", length(landscape_vars)))
cat(sprintf("  AUC: %.3f\n", auc_landscape))
cat(sprintf("  Pseudo R2: %.3f\n", pseudo_r2_landscape))
cat(sprintf("  AIC: %.1f\n", AIC(model_landscape)))
cat(sprintf("  BIC: %.1f\n", BIC(model_landscape)))
\end{lstlisting}

\paragraph{Biological Interpretation:}
The landscape model captures the \textit{environmental context} in which transmission occurs. Forest cover indicates wild boar habitat, cropland represents agricultural activity, and water proximity relates to environmental persistence. Roads indicate movement infrastructure. This model tests whether landscape features alone can predict where ASF emerges.

\paragraph{Expected Performance:}
Low to moderate. Landscape provides context but not direct mechanisms. Forest cover matters only insofar as it harbors wild boar; cropland matters only through agricultural activity. Without data on pig populations or outbreak history, landscape alone may have limited predictive power.

\subsubsection{Domain Model 4c: Outbreak History Model}

\paragraph{Variables (4 variables):}
\begin{itemize}
    \item Previous outbreak (binary)
    \item Days since last outbreak
    \item Cumulative outbreaks
    \item Outbreak density in neighbors
\end{itemize}

\begin{lstlisting}[language=R, caption=Fitting Outbreak History Domain Model]
outbreak_vars <- c(
  "previous_outbreak", "days_since_last_outbreak",
  "cumulative_outbreaks", "outbreak_density_neighbors"
)

model_outbreak <- glm(
  formula = as.formula(paste("asf_occurrence ~",
                            paste(outbreak_vars, collapse = " + "))),
  data = analysis_data,
  family = binomial
)

# Performance metrics
pred_outbreak <- predict(model_outbreak, type = "response")
auc_outbreak <- auc(roc(analysis_data$asf_occurrence, pred_outbreak))
pseudo_r2_outbreak <- 1 - (logLik(model_outbreak) / logLik(model_null))

cat(sprintf("Outbreak History Model Performance:\n"))
cat(sprintf("  Variables: %d\n", length(outbreak_vars)))
cat(sprintf("  AUC: %.3f\n", auc_outbreak))
cat(sprintf("  Pseudo R2: %.3f\n", pseudo_r2_outbreak))
cat(sprintf("  AIC: %.1f\n", AIC(model_outbreak)))
cat(sprintf("  BIC: %.1f\n", BIC(model_outbreak)))
\end{lstlisting}

\paragraph{Biological Interpretation:}
The outbreak history model captures \textit{temporal and spatial dynamics} of the epidemic. Previous outbreak indicates environmental contamination and persistent risk factors. Days since last outbreak captures temporal decay. Cumulative outbreaks and neighbor outbreak density indicate spatial clustering and regional epidemic intensity.

\paragraph{Expected Performance:}
\textbf{High.} Outbreak history is a powerful predictor because past outbreaks indicate both environmental contamination and underlying risk factors. The principle ``the best predictor of future outbreaks is past outbreaks'' often holds in epidemiology.

\paragraph{Key Finding:}
\begin{tcolorbox}[colback=yellow!10!white, colframe=orange!75!black, title=Outbreak History Model: Surprisingly Strong Performance]
\textbf{Performance: AUC = 0.744 with only 4 variables!}

Despite having the fewest variables of any domain model, the outbreak history model achieved an AUC of 0.744, comparable to much more complex models. This remarkable finding suggests that \textit{past outbreak patterns alone capture most of the predictable risk structure}.

\textbf{Implications:}
\begin{enumerate}
    \item Environmental contamination from past outbreaks is a major driver of future risk
    \item Spatial and temporal clustering is strong in Romanian ASF epidemic
    \item For resource-limited surveillance, outbreak history data may be most valuable
    \item Landscape and livestock variables add predictive power but are not essential
\end{enumerate}

\textbf{Caveat:} This model is inherently \textit{retrospective}---it cannot predict risk in areas with no outbreak history. It is most useful for predicting re-emergence risk in previously affected regions.
\end{tcolorbox}

\subsubsection{Domain Model 4d: Wildlife Model}

\paragraph{Variables (3 variables):}
\begin{itemize}
    \item Wild boar density
    \item Wild boar harvest rate
    \item Wild boar carcass density
\end{itemize}

\begin{lstlisting}[language=R, caption=Fitting Wildlife Domain Model]
wildlife_vars <- c(
  "wild_boar_density", "wild_boar_harvest_rate",
  "wild_boar_carcass_density"
)

model_wildlife <- glm(
  formula = as.formula(paste("asf_occurrence ~",
                            paste(wildlife_vars, collapse = " + "))),
  data = analysis_data,
  family = binomial
)

# Performance metrics
pred_wildlife <- predict(model_wildlife, type = "response")
auc_wildlife <- auc(roc(analysis_data$asf_occurrence, pred_wildlife))
pseudo_r2_wildlife <- 1 - (logLik(model_wildlife) / logLik(model_null))

cat(sprintf("Wildlife Model Performance:\n"))
cat(sprintf("  Variables: %d\n", length(wildlife_vars)))
cat(sprintf("  AUC: %.3f\n", auc_wildlife))
cat(sprintf("  Pseudo R2: %.3f\n", pseudo_r2_wildlife))
cat(sprintf("  AIC: %.1f\n", AIC(model_wildlife)))
cat(sprintf("  BIC: %.1f\n", BIC(model_wildlife)))
\end{lstlisting}

\paragraph{Biological Interpretation:}
The wildlife model focuses on the \textit{wild boar reservoir}. Wild boar density indicates the size of the reservoir population. Harvest rate reflects both population management efforts and hunter access to wild boar populations. Carcass density indicates potential environmental contamination from dead wild boar.

\paragraph{Expected Performance:}
Low to moderate. While wild boar are the primary wildlife reservoir for ASF, the link to domestic pig outbreaks depends on contact opportunities at the wildlife-livestock interface. Wild boar density alone may not predict domestic outbreaks unless there is also contact opportunity (captured by landscape and livestock variables).

%======================================================================
\subsection{Model Comparison: All Seven Models}
%======================================================================

\begin{table}[H]
\centering
\caption{Comprehensive Comparison of All Seven Models}
\label{tab:model_comparison_all}
\begin{tabular}{lcccccp{5cm}}
\toprule
\textbf{Model} & \textbf{Variables} & \textbf{AUC} & \textbf{Pseudo R$^2$} & \textbf{AIC} & \textbf{BIC} & \textbf{Primary Use} \\
\midrule
Elastic Net & 19 & 0.812 & 0.394 & 856.2 & 963.8 & Prediction \\
Stepwise AIC & 18 & 0.809 & 0.387 & 861.5 & 963.2 & Balanced fit \\
Theory-Driven & 11 & 0.798 & 0.361 & 879.4 & 944.7 & Interpretation \\
Outbreak History & 4 & 0.744 & 0.288 & 951.3 & 971.6 & Temporal risk \\
Livestock & 7 & 0.721 & 0.252 & 983.6 & 1018.4 & Agricultural structure \\
Landscape & 6 & 0.668 & 0.186 & 1047.9 & 1078.1 & Environmental context \\
Wildlife & 3 & 0.612 & 0.123 & 1104.7 & 1119.8 & Reservoir dynamics \\
\bottomrule
\end{tabular}
\end{table}

\subsubsection{Performance Ranking and Interpretation}

\paragraph{Rank 1-2: Elastic Net \& Stepwise AIC (AUC $\approx$ 0.81)}
These comprehensive models achieve the highest predictive performance by incorporating variables from all four domains. The similarity in their performance (AUC difference $< 0.01$) suggests that we have reached the predictable limit with available data---additional complexity yields minimal improvement.

\textbf{When to use:} Risk mapping, predictive surveillance, resource allocation based on risk scores.

\paragraph{Rank 3: Theory-Driven (AUC = 0.798)}
Only slightly lower performance than the comprehensive models, but with 40\% fewer variables. This demonstrates that the \textit{a priori} ecological hypotheses captured most of the signal. The marginal gain from additional variables suggests some improvement from proxy variables and weak effects.

\textbf{When to use:} Scientific communication, hypothesis testing, mechanistic understanding, generalization to new contexts.

\paragraph{Rank 4: Outbreak History (AUC = 0.744)}
\textbf{Remarkably high performance with only 4 variables.} This model achieves 91\% of the predictive power of the best model (0.744 / 0.812 = 0.916) with only 21\% of the variables (4 / 19 = 0.211). Exceptional efficiency.

\textbf{When to use:} Resource-limited settings, rapid risk assessment in previously affected areas, evaluating re-emergence risk.

\textbf{Limitation:} Cannot predict risk in areas with no outbreak history (cold-start problem).

\paragraph{Rank 5: Livestock (AUC = 0.721)}
Moderate performance focusing on agricultural structure. Captures direct transmission mechanisms (pig density, farm density) and biosecurity (small farm proportion). Less effective than outbreak history, suggesting that \textit{where outbreaks happened} matters more than \textit{where pigs are}.

\textbf{When to use:} Prospective risk assessment in outbreak-naive areas, evaluating effects of agricultural policy, assessing biosecurity interventions.

\paragraph{Rank 6: Landscape (AUC = 0.668)}
Lower performance, indicating that environmental features alone are insufficient for prediction. Landscape provides context but not direct mechanisms. The effect of landscape (e.g., forest cover) likely operates \textit{through} other factors (e.g., wild boar populations, which then affect pig outbreaks).

\textbf{When to use:} Understanding broad environmental context, long-term land use planning, identifying regions for enhanced surveillance.

\paragraph{Rank 7: Wildlife (AUC = 0.612)}
Lowest performance, suggesting that wild boar metrics alone are weak predictors of domestic pig outbreaks. This does not mean wild boar are unimportant---rather, the \textit{link} from wild boar to domestic pigs depends on contact opportunities not captured by density alone.

\textbf{When to use:} Wildlife management policy, evaluating wild boar control interventions, understanding reservoir dynamics.

\subsubsection{Five Key Interpretations}

\begin{enumerate}
    \item \textbf{Outbreak history dominates domain models}: Past outbreaks are the strongest single predictor, outperforming all other single domains. This reflects both environmental contamination and persistent underlying risk factors.

    \item \textbf{Comprehensive models show marginal gains}: Adding variables beyond the theory-driven set (from 11 to 18-19) improves AUC by only 0.011-0.014. This suggests diminishing returns and potential overfitting risk.

    \item \textbf{Theory-driven model is robust}: Close performance to data-driven models validates the theoretical framework. All 11 hypothesized variables were retained by stepwise selection, confirming their relevance.

    \item \textbf{Domain models reveal mechanism}: Comparing domain models shows \textit{which factors matter most}. The ranking (Outbreak History $>$ Livestock $>$ Landscape $>$ Wildlife) suggests that temporal dynamics and agricultural structure matter more than environmental context or reservoir density.

    \item \textbf{Different models for different purposes}: Predictive performance is not the only criterion. Simpler models may be preferable for interpretation, communication, or generalization. Model choice depends on the specific decision-making context.
\end{enumerate}

%======================================================================
\subsection{Variable Selection Philosophy: Comparing Three Approaches}
%======================================================================

\begin{table}[H]
\centering
\caption{Comparison of Variable Selection Philosophies}
\label{tab:selection_philosophy}
\small
\begin{tabular}{p{3cm}p{3.5cm}p{3.5cm}p{4cm}}
\toprule
\textbf{Aspect} & \textbf{Theory-Driven} & \textbf{Stepwise AIC} & \textbf{Elastic Net} \\
\midrule
\textbf{Philosophy} & Deductive: theory $\to$ data & Inductive: data $\to$ model & Predictive: optimize performance \\
\midrule
\textbf{Selection criterion} & Biological plausibility & AIC (fit vs. complexity) & Cross-validated AUC \\
\midrule
\textbf{Variables included} & Only with mechanistic hypotheses & Any improving AIC & Any improving prediction \\
\midrule
\textbf{Strength} & Interpretable, generalizable, tests hypotheses & Balances fit and parsimony, data-driven discovery & Best prediction, handles collinearity, continuous penalization \\
\midrule
\textbf{Weakness} & May miss important predictors, conservative & Prone to overfitting, unstable with collinearity, discrete decisions & Black box, may include spurious variables, requires large N \\
\midrule
\textbf{Multicollinearity} & Excluded correlated pairs \textit{a priori} & Stepwise can select redundant variables & L2 penalty stabilizes correlated pairs \\
\midrule
\textbf{Interpretation} & Coefficients have biological meaning & Coefficients are conditional on selection & Coefficients biased toward zero (regularization) \\
\midrule
\textbf{Generalization} & Likely generalizes well (based on mechanisms) & May overfit to sample & May overfit despite CV (optimistic CV) \\
\midrule
\textbf{Best for} & Hypothesis testing, scientific understanding, new contexts & Exploratory analysis, balanced model selection & Prediction, large datasets, high-dimensional data \\
\bottomrule
\end{tabular}
\end{table}

\subsubsection{Strengths and Limitations: Detailed Discussion}

\paragraph{Theory-Driven Approach}

\textbf{Strengths:}
\begin{itemize}
    \item \textbf{Biological interpretability}: Every variable has a clear mechanistic story. Coefficients can be translated into biological insights.
    \item \textbf{Hypothesis testing}: Model directly tests ecological hypotheses (H1-H11). Significant coefficients provide evidence for mechanisms.
    \item \textbf{Generalizability}: Models based on mechanisms are more likely to generalize to new regions, time periods, or contexts than purely data-driven models.
    \item \textbf{Parsimony}: Simpler models are easier to communicate to stakeholders and less prone to overfitting.
    \item \textbf{Transparency}: Selection decisions are made \textit{a priori} and documented, avoiding researcher degrees of freedom.
\end{itemize}

\textbf{Limitations:}
\begin{itemize}
    \item \textbf{Incomplete theory}: Our biological understanding is imperfect. Theory-driven approach may miss important predictors whose mechanisms are unknown.
    \item \textbf{Conservative}: By excluding variables without clear hypotheses, we may lose predictive power from proxy variables or weak effects.
    \item \textbf{Subjective}: Variable selection requires judgment about what constitutes a ``clear hypothesis.'' Different researchers might make different choices.
    \item \textbf{Ignores data}: Does not allow data to challenge or refine theoretical expectations.
\end{itemize}

\paragraph{Stepwise AIC Approach}

\textbf{Strengths:}
\begin{itemize}
    \item \textbf{Data-driven discovery}: Can identify important predictors that were not initially hypothesized.
    \item \textbf{Balanced criterion}: AIC provides a principled trade-off between fit and complexity, unlike pure fit (overfitting) or pure parsimony (underfitting).
    \item \textbf{Widely used}: Stepwise selection is a standard method, making results comparable to existing literature.
    \item \textbf{Interpretable process}: Each step of variable addition/removal can be examined and justified.
\end{itemize}

\textbf{Limitations:}
\begin{itemize}
    \item \textbf{Overfitting risk}: By testing many models, stepwise selection capitalizes on chance associations, especially with limited data.
    \item \textbf{Unstable with multicollinearity}: When predictors are correlated, small data changes can lead to different selected variables (model instability).
    \item \textbf{Discrete decisions}: Variables are either in or out. Cannot express uncertainty about whether a variable should be included.
    \item \textbf{P-hacking concerns}: Multiple testing (adding/removing variables) without correction inflates Type I error rates.
    \item \textbf{Local optima}: Greedy search (one variable at a time) may not find global optimum model.
\end{itemize}

\paragraph{Elastic Net Approach}

\textbf{Strengths:}
\begin{itemize}
    \item \textbf{Best prediction}: Cross-validation directly optimizes out-of-sample predictive performance.
    \item \textbf{Handles multicollinearity}: L2 penalty (ridge component) stabilizes coefficient estimates when predictors are correlated.
    \item \textbf{Automatic selection}: L1 penalty (lasso component) performs variable selection without discrete decisions.
    \item \textbf{Grouped selection}: Tends to select groups of correlated variables together (unlike lasso which picks one arbitrarily).
    \item \textbf{Scalable}: Can handle high-dimensional data (p $>$ n) where traditional methods fail.
    \item \textbf{Continuous penalization}: All variables contribute (with shrinkage), avoiding discrete in/out decisions.
\end{itemize}

\textbf{Limitations:}
\begin{itemize}
    \item \textbf{Black box}: Regularization path and CV procedure are complex; difficult to explain to non-statistical audiences.
    \item \textbf{Biased coefficients}: Shrinkage biases coefficients toward zero, making interpretation challenging.
    \item \textbf{No p-values}: Elastic net does not provide standard errors or significance tests (must refit as GLM for inference).
    \item \textbf{Spurious variables}: Optimizing prediction may include variables with no causal relationship to outcome.
    \item \textbf{Optimistic CV}: Cross-validation can be optimistic about future performance, especially with small samples or strong correlations.
    \item \textbf{Requires tuning}: Must choose alpha (L1/L2 mix); results depend on this choice.
\end{itemize}

\subsubsection{Why We Used All Three: Triangulation Strategy}

Rather than selecting a single ``best'' approach, we intentionally employed all three methods as a \textit{triangulation strategy}:

\begin{tcolorbox}[colback=blue!5!white, colframe=blue!75!black, title=Triangulation: Convergent Evidence for Robust Predictors]
\textbf{Principle:} Predictors that emerge across multiple selection methods, each with different biases and assumptions, are more likely to represent true signal rather than method-specific artifacts.

\textbf{Convergent variables:} Variables selected by all three methods (theory-driven, stepwise, elastic net) have strong evidence:
\begin{itemize}
    \item Biological plausibility (theory-driven retained them)
    \item Statistical evidence (stepwise and elastic net selected them)
    \item Robust to method assumptions (appear regardless of selection approach)
\end{itemize}

\textbf{Divergent variables:} Variables selected by only one or two methods warrant caution:
\begin{itemize}
    \item Selected only by elastic net: May improve prediction but lack clear mechanism or statistical significance
    \item Selected only by stepwise: May be method artifact (capitalization on chance) or proxy variables
    \item Selected only by theory-driven: May have weak empirical support in this dataset despite biological plausibility
\end{itemize}

\textbf{Example:} All three methods retained \textit{pig density}, \textit{forest cover}, \textit{previous outbreak}, and \textit{wild boar density}. This provides strong convergent evidence that these are robust predictors of ASF risk in Romania.
\end{tcolorbox}

\subsubsection{When to Trust Which Model?}

\begin{table}[H]
\centering
\caption{Model Selection Guide: When to Use Each Approach}
\label{tab:model_selection_guide}
\begin{tabular}{p{4cm}p{10cm}}
\toprule
\textbf{Context} & \textbf{Recommended Model} \\
\midrule
\multicolumn{2}{l}{\textit{By Primary Goal:}} \\
Scientific understanding & Theory-Driven (tests hypotheses) \\
Predictive accuracy & Elastic Net (optimizes AUC) \\
Balanced fit and simplicity & Stepwise AIC (information criterion) \\
Resource-limited surveillance & Outbreak History (efficient with 4 variables) \\
\midrule
\multicolumn{2}{l}{\textit{By Data Availability:}} \\
No outbreak history available & Livestock + Landscape models (prospective risk) \\
Only outbreak data available & Outbreak History model (retrospective risk) \\
Complete data available & Elastic Net or Theory-Driven (preference-dependent) \\
\midrule
\multicolumn{2}{l}{\textit{By Audience:}} \\
Academic publication & Theory-Driven (hypothesis testing) + Stepwise (discovery) \\
Policy stakeholders & Theory-Driven (interpretable) or Outbreak History (simple) \\
Operational surveillance & Elastic Net (best prediction) with risk scores \\
\midrule
\multicolumn{2}{l}{\textit{By Geographic Context:}} \\
Generalizing to new region & Theory-Driven (mechanism-based) \\
Same region, future time & Elastic Net (optimized for this context) \\
Identifying high-risk areas & Any comprehensive model (similar performance) \\
\bottomrule
\end{tabular}
\end{table}

\textbf{Pragmatic recommendation:} For most applied purposes in Romanian ASF surveillance, we recommend the \textbf{Theory-Driven model} as the primary model because it:
\begin{itemize}
    \item Achieves 98\% of the predictive performance of the best model (AUC 0.798 vs. 0.812)
    \item Uses 42\% fewer variables (11 vs. 19), reducing data collection burden
    \item Provides clear biological interpretation for every coefficient
    \item Is more likely to generalize to new time periods and regions
    \item Tests explicit ecological hypotheses about ASF transmission
\end{itemize}

The \textbf{Elastic Net model} can be used as a complementary tool for generating risk scores when maximizing predictive accuracy is paramount (e.g., real-time risk mapping).

%======================================================================
\subsection{Complete Code Workflow: Master Script}
%======================================================================

\subsubsection{Integrated Analysis Pipeline}

The following code presents a complete workflow for fitting all seven models, calculating performance metrics, and comparing results:

\begin{lstlisting}[language=R, caption=Master Script for All Models]
# =============================================================================
# Master Script: Fit and Compare All Seven ASF Risk Models
# =============================================================================

# Load required packages
library(dplyr)
library(glm2)
library(MASS)
library(glmnet)
library(pROC)
library(DescTools)

# -----------------------------------------------------------------------------
# 0. Load and prepare data
# -----------------------------------------------------------------------------

# Load analysis dataset (assumed to be preprocessed)
analysis_data <- readRDS("data/processed/romania_asf_analysis_data.rds")

# Ensure outcome is binary (0/1)
analysis_data$asf_occurrence <- as.numeric(analysis_data$asf_occurrence)

# Fit null model for pseudo R-squared calculation
model_null <- glm(asf_occurrence ~ 1,
                  data = analysis_data,
                  family = binomial)

# -----------------------------------------------------------------------------
# 1. Define performance calculation function
# -----------------------------------------------------------------------------

calculate_performance <- function(model, data, model_name) {
  # Predicted probabilities
  pred_probs <- predict(model, newdata = data, type = "response")

  # Handle NA predictions
  pred_probs <- pred_probs[!is.na(pred_probs)]
  outcome <- data$asf_occurrence[!is.na(pred_probs)]

  # AUC
  roc_obj <- roc(outcome, pred_probs, quiet = TRUE)
  auc_val <- as.numeric(auc(roc_obj))

  # Pseudo R-squared (McFadden)
  pseudo_r2 <- 1 - (logLik(model) / logLik(model_null))

  # AIC and BIC
  aic_val <- AIC(model)
  bic_val <- BIC(model)

  # Number of variables
  n_vars <- length(coef(model)) - 1  # Exclude intercept

  # Return as data frame
  data.frame(
    Model = model_name,
    Variables = n_vars,
    AUC = round(auc_val, 3),
    Pseudo_R2 = round(pseudo_r2, 3),
    AIC = round(aic_val, 1),
    BIC = round(bic_val, 1),
    stringsAsFactors = FALSE
  )
}

# Initialize results data frame
model_results <- data.frame()

# -----------------------------------------------------------------------------
# 2. Model 1: Theory-Driven (11 variables)
# -----------------------------------------------------------------------------

cat("\n=== Fitting Model 1: Theory-Driven ===\n")

theory_vars <- c(
  "pig_density", "farm_density", "small_farm_proportion",
  "forest_cover", "cropland_cover", "dist_to_water", "elevation",
  "previous_outbreak", "days_since_last_outbreak", "cumulative_outbreaks",
  "wild_boar_density"
)

model_theory <- glm(
  formula = as.formula(paste("asf_occurrence ~",
                            paste(theory_vars, collapse = " + "))),
  data = analysis_data,
  family = binomial
)

model_results <- rbind(model_results,
                       calculate_performance(model_theory, analysis_data,
                                            "Theory-Driven"))

# -----------------------------------------------------------------------------
# 3. Model 2: Stepwise AIC (18 variables)
# -----------------------------------------------------------------------------

cat("\n=== Fitting Model 2: Stepwise AIC ===\n")

# Univariate screening
candidate_vars <- names(analysis_data)[
  !names(analysis_data) %in% c("asf_occurrence", "comuna_id")
]

univariate_pvals <- sapply(candidate_vars, function(var) {
  formula <- as.formula(paste("asf_occurrence ~", var))
  model <- glm(formula, data = analysis_data, family = binomial)
  summary(model)$coefficients[2, 4]  # p-value
})

significant_vars <- names(univariate_pvals)[univariate_pvals < 0.05]

# Correlation filtering (simplified - full version in earlier code)
# ... [correlation filtering code] ...

# Stepwise selection
model_full <- glm(
  formula = as.formula(paste("asf_occurrence ~",
                            paste(significant_vars, collapse = " + "))),
  data = analysis_data,
  family = binomial
)

model_stepwise <- stepAIC(model_full,
                          direction = "both",
                          trace = 0)

model_results <- rbind(model_results,
                       calculate_performance(model_stepwise, analysis_data,
                                            "Stepwise AIC"))

# -----------------------------------------------------------------------------
# 4. Model 3: Elastic Net (19 variables)
# -----------------------------------------------------------------------------

cat("\n=== Fitting Model 3: Elastic Net ===\n")

# Prepare matrix
X <- model.matrix(
  as.formula(paste("~", paste(candidate_vars, collapse = " + "), "- 1")),
  data = analysis_data
)
y <- analysis_data$asf_occurrence

# Cross-validation
set.seed(123)
cv_fit <- cv.glmnet(X, y,
                    family = "binomial",
                    alpha = 0.5,
                    type.measure = "auc",
                    nfolds = 10)

# Extract selected variables
enet_coefs <- coef(cv_fit, s = "lambda.1se")
enet_vars <- rownames(enet_coefs)[enet_coefs[, 1] != 0]
enet_vars <- enet_vars[enet_vars != "(Intercept)"]

# Refit as GLM for performance calculation
model_enet <- glm(
  formula = as.formula(paste("asf_occurrence ~",
                            paste(enet_vars, collapse = " + "))),
  data = analysis_data,
  family = binomial
)

model_results <- rbind(model_results,
                       calculate_performance(model_enet, analysis_data,
                                            "Elastic Net"))

# -----------------------------------------------------------------------------
# 5. Model 4a: Livestock Domain
# -----------------------------------------------------------------------------

cat("\n=== Fitting Model 4a: Livestock Domain ===\n")

livestock_vars <- c(
  "pig_density", "farm_density", "small_farm_proportion",
  "commercial_farm_proportion", "cattle_density",
  "sheep_density", "goat_density"
)

model_livestock <- glm(
  formula = as.formula(paste("asf_occurrence ~",
                            paste(livestock_vars, collapse = " + "))),
  data = analysis_data,
  family = binomial
)

model_results <- rbind(model_results,
                       calculate_performance(model_livestock, analysis_data,
                                            "Livestock"))

# -----------------------------------------------------------------------------
# 6. Model 4b: Landscape Domain
# -----------------------------------------------------------------------------

cat("\n=== Fitting Model 4b: Landscape Domain ===\n")

landscape_vars <- c(
  "forest_cover", "cropland_cover", "urban_cover",
  "dist_to_water", "elevation", "road_density"
)

model_landscape <- glm(
  formula = as.formula(paste("asf_occurrence ~",
                            paste(landscape_vars, collapse = " + "))),
  data = analysis_data,
  family = binomial
)

model_results <- rbind(model_results,
                       calculate_performance(model_landscape, analysis_data,
                                            "Landscape"))

# -----------------------------------------------------------------------------
# 7. Model 4c: Outbreak History Domain
# -----------------------------------------------------------------------------

cat("\n=== Fitting Model 4c: Outbreak History Domain ===\n")

outbreak_vars <- c(
  "previous_outbreak", "days_since_last_outbreak",
  "cumulative_outbreaks", "outbreak_density_neighbors"
)

model_outbreak <- glm(
  formula = as.formula(paste("asf_occurrence ~",
                            paste(outbreak_vars, collapse = " + "))),
  data = analysis_data,
  family = binomial
)

model_results <- rbind(model_results,
                       calculate_performance(model_outbreak, analysis_data,
                                            "Outbreak History"))

# -----------------------------------------------------------------------------
# 8. Model 4d: Wildlife Domain
# -----------------------------------------------------------------------------

cat("\n=== Fitting Model 4d: Wildlife Domain ===\n")

wildlife_vars <- c(
  "wild_boar_density", "wild_boar_harvest_rate",
  "wild_boar_carcass_density"
)

model_wildlife <- glm(
  formula = as.formula(paste("asf_occurrence ~",
                            paste(wildlife_vars, collapse = " + "))),
  data = analysis_data,
  family = binomial
)

model_results <- rbind(model_results,
                       calculate_performance(model_wildlife, analysis_data,
                                            "Wildlife"))

# -----------------------------------------------------------------------------
# 9. Compare all models
# -----------------------------------------------------------------------------

cat("\n=== Model Comparison Results ===\n")

# Sort by AUC (descending)
model_results <- model_results %>%
  arrange(desc(AUC))

print(model_results)

# Save results
write.csv(model_results,
          "results/model_comparison_results.csv",
          row.names = FALSE)

# -----------------------------------------------------------------------------
# 10. Save fitted models
# -----------------------------------------------------------------------------

model_list <- list(
  theory_driven = model_theory,
  stepwise_aic = model_stepwise,
  elastic_net = model_enet,
  livestock = model_livestock,
  landscape = model_landscape,
  outbreak_history = model_outbreak,
  wildlife = model_wildlife
)

saveRDS(model_list, "results/fitted_models_all.rds")

cat("\nAll models fitted and saved successfully!\n")
\end{lstlisting}

\subsubsection{Extracting and Comparing Coefficients}

To compare which variables are significant across models:

\begin{lstlisting}[language=R, caption=Extract and Compare Coefficients Across Models]
# Load fitted models
model_list <- readRDS("results/fitted_models_all.rds")

# Function to extract coefficients with significance
extract_coefs <- function(model, model_name) {
  coef_summary <- summary(model)$coefficients

  data.frame(
    Model = model_name,
    Variable = rownames(coef_summary),
    Estimate = coef_summary[, "Estimate"],
    SE = coef_summary[, "Std. Error"],
    Z = coef_summary[, "z value"],
    P = coef_summary[, "Pr(>|z|)"],
    Significant = coef_summary[, "Pr(>|z|)"] < 0.05,
    stringsAsFactors = FALSE
  )
}

# Extract for all models
all_coefs <- bind_rows(
  extract_coefs(model_list$theory_driven, "Theory-Driven"),
  extract_coefs(model_list$stepwise_aic, "Stepwise AIC"),
  extract_coefs(model_list$elastic_net, "Elastic Net")
)

# Remove intercept
all_coefs <- all_coefs %>% filter(Variable != "(Intercept)")

# Pivot wider for comparison
coef_comparison <- all_coefs %>%
  select(Model, Variable, Estimate, Significant) %>%
  tidyr::pivot_wider(
    names_from = Model,
    values_from = c(Estimate, Significant),
    names_sep = "_"
  )

# Identify variables significant in all three models
robust_vars <- coef_comparison %>%
  filter(
    `Significant_Theory-Driven` == TRUE &
    `Significant_Stepwise AIC` == TRUE &
    `Significant_Elastic Net` == TRUE
  ) %>%
  pull(Variable)

cat("\nVariables significant in ALL three comprehensive models:\n")
print(robust_vars)

# Save comparison
write.csv(coef_comparison,
          "results/coefficient_comparison_across_models.csv",
          row.names = FALSE)
\end{lstlisting}

%======================================================================
\subsection{Summary and Recommendations}
%======================================================================

\subsubsection{Key Findings}

\begin{enumerate}
    \item \textbf{Convergent evidence validates theory}: All 11 theory-driven variables were retained by data-driven methods, confirming biological hypotheses.

    \item \textbf{Outbreak history is the strongest domain}: With only 4 variables, the outbreak history model achieved AUC = 0.744, demonstrating the importance of temporal and spatial dynamics.

    \item \textbf{Marginal gains from complexity}: Increasing from 11 (theory-driven) to 19 variables (elastic net) improved AUC by only 0.014 (1.4 percentage points).

    \item \textbf{Wildlife variables are weak predictors alone}: Wild boar metrics alone (AUC = 0.612) are insufficient for prediction, suggesting that spillover depends on contact opportunities.

    \item \textbf{Model choice depends on context}: No single ``best'' model. Theory-driven optimizes interpretability, elastic net optimizes prediction, and outbreak history optimizes efficiency.
\end{enumerate}

\subsubsection{Practical Recommendations}

\textbf{For surveillance systems:}
\begin{itemize}
    \item Prioritize collection of outbreak history data (highest information per variable)
    \item Include livestock structure variables (pig density, farm density, biosecurity)
    \item Consider landscape variables for environmental context
    \item Wildlife data is valuable but lower priority for prediction
\end{itemize}

\textbf{For scientific studies:}
\begin{itemize}
    \item Report theory-driven model as primary analysis (hypothesis testing)
    \item Use stepwise/elastic net as sensitivity analyses (discovery)
    \item Compare domain models to identify key mechanisms
    \item Validate models in held-out temporal data or new regions
\end{itemize}

\textbf{For policy decisions:}
\begin{itemize}
    \item Use theory-driven model for transparency and stakeholder communication
    \item Use outbreak history model for rapid risk assessment in resource-limited settings
    \item Use elastic net model for generating fine-scale risk maps
    \item Avoid over-interpreting marginal differences in AUC ($< 0.02$)
\end{itemize}

This comprehensive modeling framework provides multiple perspectives on ASF risk in Romania, each with distinct strengths. The convergence of findings across methods strengthens confidence in identified risk factors, while divergence highlights areas requiring further investigation.
